\appendix
\chapter{Sursele și proveniența figurilor}
\label{anexa:surse-figuri}
\pagestyle{fancy}

Tabelele de mai jos indică sursele și proveniența imaginilor folosite în lucrare.
Numele figurii este un link intern către figura corespunzătoare din capitolul în care este folosită.

\begin{table}[ht]
    \caption{Sursele figurilor folosite în capitolul de analiză}
    \centering
    \footnotesize
    \begin{tabular}{|p{0.58\textwidth}|p{0.30\textwidth}|}
        \hline
        Figură & Link sursă \\
        \hline
        \hyperref[fig:yoloe-architecture]{Figura~\ref*{fig:yoloe-architecture} -- Arhitectura YOLOE} &
        \href{https://docs.ultralytics.com/models/yoloe/#architecture-overview}{Sursă YOLOE} \\
        \hline
        \hyperref[fig:mask2former-architecture]{Figura~\ref*{fig:mask2former-architecture} -- Arhitectura Mask2Former} &
        \href{https://arxiv.org/pdf/2112.01527}{Sursă Mask2Former} \\
        \hline
        \hyperref[fig:transformer-decoder-only]{Figura~\ref*{fig:transformer-decoder-only} -- Transformer decoder-only} &
        \href{https://substack-post-media.s3.amazonaws.com/public/images/f6133c18-bfaf-4578-8c5a-e5ac7809f65b_1632x784.png}{Sursă Transformer} \\
        \hline
        \hyperref[fig:jwt-structure]{Figura~\ref*{fig:jwt-structure} -- Structura JWT} &
        \href{https://monsterlessons-academy.com/posts/jwt-authentication-beginners-guide-with-real-application}{Sursă JWT} \\
        \hline
    \end{tabular}
    \label{tab:surse-figuri-ai}
\end{table}

\begin{table}[ht]
    \caption{Proveniența figurilor folosite în capitolul de proiectare și implementare}
    \centering
    \footnotesize
    \begin{tabular}{|p{0.78\textwidth}|p{0.08\textwidth}|}
        \hline
        Figură & Proveniență \\
        \hline
        \hyperref[fig:arhitectura-generala-aplicatie]{Figura~\ref*{fig:arhitectura-generala-aplicatie} -- Arhitectura generală a aplicației} &
        - \\
        \hline
        \hyperref[fig:pyvisionannotator-interfata-principala]{Figura~\ref*{fig:pyvisionannotator-interfata-principala} -- Interfața principală a aplicației PyVisionAnnotator} &
        - \\
        \hline
        \hyperref[fig:organizare-mainwindow]{Figura~\ref*{fig:organizare-mainwindow} -- Organizarea internă a ferestrei principale} &
        - \\
        \hline
        \hyperref[fig:setari-autoseg-automask]{Figura~\ref*{fig:setari-autoseg-automask} -- Panoul de setări pentru AutoSeg și AutoMask} &
        - \\
        \hline
        \hyperref[fig:panou-proprietati-adnotari]{Figura~\ref*{fig:panou-proprietati-adnotari} -- Panoul de proprietăți și lista de adnotări} &
        - \\
        \hline
        \hyperref[fig:flux-evenimente-canvas]{Figura~\ref*{fig:flux-evenimente-canvas} -- Fluxul de evenimente în canvas} &
        - \\
        \hline
        \hyperref[fig:clase-principale-adnotari]{Figura~\ref*{fig:clase-principale-adnotari} -- Clasele principale pentru adnotări} &
        - \\
        \hline
        \hyperref[fig:flux-autoseg-yolo-client]{Figura~\ref*{fig:flux-autoseg-yolo-client} -- Fluxul AutoSeg YOLO în client} &
        - \\
        \hline
        \hyperref[fig:flux-automask-punct-client]{Figura~\ref*{fig:flux-automask-punct-client} -- Fluxul AutoMask pe baza unui punct selectat} &
        - \\
        \hline
        \hyperref[fig:fereastra-chat-colaborare]{Figura~\ref*{fig:fereastra-chat-colaborare} -- Fereastra de chat și colaborare} &
        - \\
        \hline
        \hyperref[fig:fereastra-chatbot-local]{Figura~\ref*{fig:fereastra-chatbot-local} -- Fereastra chatbot local cu Ollama sau llama.cpp} &
        - \\
        \hline
        \hyperref[fig:flux-incarcare-protectie-nesalvat]{Figura~\ref*{fig:flux-incarcare-protectie-nesalvat} -- Fluxul de încărcare imagine și protecția modificărilor nesalvate} &
        - \\
        \hline
        \hyperref[fig:flux-operational-colaborare]{Figura~\ref*{fig:flux-operational-colaborare} -- Fluxul operațional al colaborării dintre utilizator, clientul desktop și backend} &
        - \\
        \hline
    \end{tabular}
    \label{tab:provenienta-figuri-implementare}
\end{table}
